% !TITLE 送友人
% !AUTHOR 李白
% !DYNASTY 唐
% !ORDER 12

\begin{poem}
青山横北郭\textsuperscript{1}，白水绕东城。
此地一为别，孤蓬\textsuperscript{2}万里征。
浮云游子意，落日故人情。
挥手自兹去，萧萧班马鸣\textsuperscript{3}。
\end{poem}

\begin{annotations}
\notetext{1}{郭：外城。青山横在北城之外，白水环绕着东城。}
\notetext{2}{孤蓬：蓬草干枯后根断，随风飘转，比喻漂泊无定的人。你就要像孤独的蓬草一样远行万里了。}
\notetext{3}{萧萧班马鸣：离别时马发出萧萧的悲鸣。班马，离群的马。据说马通人性，当骑手分别时马也会感到悲伤。}
\end{annotations}

\begin{appreciation}
这是李白写得最工整、最含蓄的送别诗之一。李白这个人，高兴时写诗不管格律，可一到送别，规矩得像换了个人。

青山横北郭，白水绕东城。郭是外城。城北横着一脉青山，城东绕着一道白水。两句对得极工整：青对白，山对水，横对绕，北对东，颜色、动静、方位，一样一样全配好了。送别的地点，就定在这幅画里。

此地一为别，孤蓬万里征。蓬是蓬草，枯了以后根就断，风一吹满地滚，古诗里专门用来比喻漂泊的人。此地一别，你就要像一株断了根的蓬草，滚到万里之外去。一个“孤蓬”，把这趟远行的孤单全说完了。

浮云游子意，落日故人情。这一联是全诗的眼睛。天上飘着浮云，浮云没有定所，游子也是；太阳落得那么慢，舍不得走完，故人也是。把感情放进眼前的景物里，景物就替人说了话。整首诗没有一个字直说舍不得，但你读到这里就全懂了。

挥手自兹去，萧萧班马鸣。班马是离群的马。人已挥完手，该说的都说完了，只剩两匹马隔着路互相嘶鸣。马都懂的离别，人反而说不出口——这一声马鸣，比任何一句“珍重”都动人。

我们读过王勃的《送杜少府之任蜀州》，少年王勃说“海内存知己，天涯若比邻”，气概大得很。李白不那样写。写这首送别诗时，李白已是中年人，见过太多“此地一为别”之后再也不见的故人，所以他把话都咽了回去，只写了青山、浮云、落日和马鸣。年轻时的送别是豪言，中年的送别是静默，没有谁更高明，只是人到了一定的年纪，都这样。

将来你去远方读书，爸妈去车站送你，大概也是这个场面——检票口一分开，话就只剩“路上小心”。古人送别的叮嘱，哥哥也记着那句：“努力加餐饭”。不管走到哪，好好吃饭，家里人最惦记的就是这个。
\end{appreciation}
